\documentclass[10pt,a4paper]{article}
\usepackage[margin=2.1cm]{geometry}
\usepackage{fontspec}
\usepackage{xeCJK}
\setCJKmainfont{Noto Serif CJK KR}
\setCJKsansfont{Noto Sans CJK KR}
\setCJKmonofont{Noto Sans Mono CJK KR}
\xeCJKsetup{CJKspace=true}
\usepackage{amsmath,amssymb,mathtools}
\usepackage{booktabs,tabularx,multirow,array,longtable}
\usepackage{enumitem}
\usepackage{xcolor}
\usepackage{caption}
\usepackage{tikz}
\usetikzlibrary{positioning,arrows.meta,fit,backgrounds,calc,shapes.geometric,decorations.pathreplacing}
\usepackage[colorlinks=true]{hyperref}
\linespread{1.15}
\setlength{\parskip}{3pt}
\setlist{itemsep=1pt,topsep=2pt}

% ---------- notation macros ----------
\newcommand{\G}{\mathcal{G}}
\newcommand{\KS}{\mathcal{K}_S}
\newcommand{\Ont}{\mathcal{O}}
\newcommand{\Ctx}{\mathcal{C}}
\newcommand{\Tcmd}{\mathcal{T}_{\mathrm{cmd}}}
\newcommand{\Treq}{\mathcal{T}_{\mathrm{req}}}
\newcommand{\Tst}{\mathcal{T}_{\mathrm{state}}}
\newcommand{\Tadv}{\mathcal{T}_{\mathrm{adv}}}
\newcommand{\Yset}{\mathcal{Y}}
\newcommand{\Vdom}{\mathcal{V}_{\mathrm{dom}}}
\newcommand{\KR}{\mathcal{K}_R}
\newcommand{\Cfg}{\Theta_{\mathrm{kg}}}
\newcommand{\Ivl}{\mathbb{I}}
\newcommand{\ATCO}{\textsf{ATCO}}
\newcommand{\PILOT}{\textsf{PILOT}}
\newcommand{\cmd}[1]{\textsf{\small #1}}
\newcommand{\ch}[1]{\textbf{\textcolor{red!60!black}{#1}}}
\newcommand{\md}[1]{\textbf{\textcolor{blue!60!black}{#1}}}

\definecolor{perc}{RGB}{225,238,250}
\definecolor{comp}{RGB}{232,245,228}
\definecolor{proj}{RGB}{252,240,222}
\definecolor{res}{RGB}{240,236,248}
\pagestyle{empty}
\begin{document}
\section*{KG 스키마와 시간 모델 (최종)}
노드를 성격에 따라 세 층으로 나눈다: \textbf{① 사건 층}(시각을 가짐), \textbf{② 개체 층}(송신에서 생기지만 시각 없음), \textbf{③ 정적 KB 층}(미리 준비한 지식). 아래 표는 이 그림이 담은 설계 결정과 근거이고, 맨 아래는 KATL 실제 교신으로 그린 시간축 예시다.
\begin{center}
\resizebox{0.92\textwidth}{!}{%
\begin{tikzpicture}[
  font=\small,
  dyn/.style={draw, thick, rounded corners=2pt, fill=comp, minimum height=0.85cm, minimum width=2.6cm, align=center},
  ent/.style={draw, thick, rounded corners=2pt, fill=yellow!12, minimum height=0.85cm, minimum width=2.6cm, align=center},
  sta/.style={draw, thick, fill=res, minimum height=0.85cm, minimum width=2.4cm, align=center},
  e/.style={-{Stealth[length=2.2mm]}, thick},
  el/.style={font=\scriptsize\sffamily, fill=white, inner sep=1pt},
  band/.style={draw, rounded corners=6pt, thick, inner sep=0pt},
  bl/.style={font=\small\bfseries, anchor=north west, align=left}
]
% ---- layer 1: events (time-stamped)
\node[dyn] (tx) at (0,0) {\textbf{Transmission}\\{\scriptsize $s=[t^s,t^e)$, $w$, speaker}};
\node[dyn] (ins) at (7,0) {\textbf{Instruction}\\{\scriptsize $\kappa,v,q$, $\tau=[t^s,t^{\mathrm{to}})$}};
\draw[e] (tx) -- node[el,above]{CONTAINS} (ins);
\draw[e, red!60!black] (tx.north west) to[out=110,in=180,looseness=1.2] ($(tx.north)+(-0.2,1.0)$) node[el,above,text=red!60!black,align=center]{READBACK\_OF \{status\}\\{\scriptsize PILOT 송신 $\to$ 직전 ATCO 송신, 같은 항공편}} to[out=0,in=70,looseness=1.2] (tx.north);
\draw[e, teal!60!black] (ins.north west) to[out=110,in=180,looseness=1.2] ($(ins.north)+(-0.2,1.0)$) node[el,above,text=teal!60!black,align=center]{NEXT \{$\Delta t$\}\\{\scriptsize 같은 항공편의 다음 관제사 지시}} to[out=0,in=70,looseness=1.2] (ins.north);
\draw[-{Stealth[length=2.5mm]}, very thick, teal!60!black] (-3.2,-1.1) -- node[below,font=\scriptsize]{시간 $t$ (UTC)} (10.6,-1.1);
% ---- layer 2: entities (from transmissions, no time)
\node[ent] (fl) at (0,-3.6) {\textbf{Flight}\\{\scriptsize $c$ (예: DAL1489)}};
\node[ent] (cm) at (-4.2,-3.6) {\textbf{CallsignMention}\\{\scriptsize $\hat c$ (원 표면형)}};
% ---- layer 3: static KB
\node[sta] (role) at (-4.2,-6.6) {\textbf{Role}\\{\scriptsize ATCO / PILOT}};
\node[sta] (al) at (0,-6.6) {\textbf{Airline}\\{\scriptsize designator, country (예: DAL, US)}};
\node[sta, minimum width=3.6cm] (rw) at (7,-6.6) {\textbf{지시 대상}\\{\scriptsize Runway $\cdot$ Taxiway $\cdot$ Ramp $\cdot$ Facility}\\{\scriptsize 예: \texttt{Taxiway:KATL:M2}}};
\node[sta] (ap) at (12.9,-6.6) {\textbf{Airport}\\{\scriptsize location\_indicator,}\\{\scriptsize country (예: KATL, US)}};
% ---- edges across layers
\draw[e] (tx) -- node[el,right]{ADDRESSES} (fl);
\draw[e] (fl) -- node[el,above]{ALIAS} (cm);
\draw[e] (fl) -- node[el,right]{OPERATED\_BY} (al);
\draw[e] (tx.south west) to[out=225,in=60] node[el,pos=0.35,left]{SPOKEN\_BY} (role.north);
\draw[e] (ins.south) -- node[el,right,pos=0.55,align=left]{TARGET or VIA\\{\scriptsize VIA는 Taxiway만 (taxi 경로)}} (rw.north);
\draw[e] (rw) -- node[el,above]{PART\_OF} (ap);
% ---- bands
\begin{scope}[on background layer]
\node[band, draw=teal!60!black, fill=teal!5, fit={(-10,2.6) (14.8,-1.6)}] (b1) {};
\node[band, draw=orange!70!black, fill=yellow!6, fit={(-10,-2.6) (14.8,-4.5)}] (b2) {};
\node[band, draw=violet!60!black, fill=violet!4, fit={(-10,-5.6) (14.8,-7.5)}] (b3) {};
\end{scope}
\node[bl, text=teal!60!black] at ([xshift=2mm,yshift=-2mm]b1.north west) {① 사건 층\\{\scriptsize\mdseries 시각을 가짐}\\{\scriptsize\mdseries (Zep: episode)}};
\node[bl, text=orange!70!black] at ([xshift=2mm,yshift=-2mm]b2.north west) {② 개체 층\\{\scriptsize\mdseries 송신에서 생김,}\\{\scriptsize\mdseries 시각 없음}};
\node[bl, text=violet!60!black] at ([xshift=2mm,yshift=-2mm]b3.north west) {③ 정적 KB 층\\{\scriptsize\mdseries 미리 준비한 지식}};
\end{tikzpicture}}
\end{center}
\begin{center}\footnotesize
\begin{tabularx}{\textwidth}{@{}>{\raggedright\arraybackslash}p{3.0cm}>{\raggedright\arraybackslash}X>{\raggedright\arraybackslash}p{3.6cm}@{}}
\toprule
결정 & 내용 & 근거\\\midrule
세 층으로 나눔 & ① 사건 층(Transmission, Instruction: 시각 있음) ② 개체 층(Flight, CallsignMention: 송신에서 생기나 시각 없음) ③ 정적 KB 층 & Zep의 episode / entity 층 구분\\
송신·지시를 노드로 (사건화) & 송신 하나에 지시가 여러 개(DLR: Callsign + Command 여러 개)라 두 노드를 합치지 않음 & W3C n-ary 관계, SEM 사건 노드\\
시간은 속성으로 & 송신 $s=[t^s,t^e)$, 지시 $\tau=[t^s,t^{\mathrm{to}})$, 간선 적재 시각 $t'$. 시각이 연속값이라 시간 노드를 공유할 일이 없음. \texttt{valid\_from/valid\_to}는 OWL-Time \texttt{hasBeginning/hasEnd}에 대응 & SEM \texttt{hasBeginTimeStamp}, Zep 이중 시간($T$, $T'$)\\
시점 질의 & 시각 $t$를 주면 $\tau\ni t$인 관제사 상태형 지시를 고르고(\texttt{snapshot}), \texttt{tx\_node}로 송신을 찾음 & \\
NEXT는 지시끼리 & 항공기 상태의 흐름(어떤 지시를 어떤 순서로 받았나). 한 송신 안 지시 순서도 담고(\,$\Delta t=0$), 이전 지시를 끝내는 것도 지시 단위. 시작 시각 순서에서 유도됨 & 사건 그래프의 next-event 관계\\
READBACK\_OF는 송신끼리 & 대화의 주고받기(누가 누구에게 답했나). 조건: PILOT $\to$ 직전 ATCO, 같은 항공편·주파수(복창 시간 창은 \texttt{kg\_build.json}). 빠뜨린 지시는 상대 노드가 없어 지시 간선으로 못 그리므로, 지시 단위 결과는 간선 속성(\texttt{matched}, \texttt{missing}, \texttt{conflicts})에 저장 & 대화 구조 연구의 reply-to\\
지우지 않고 끝냄 & 새 지시가 이전 지시를 대체하면 $\tau$의 끝과 \texttt{invalidated\_by}만 기록. 과거 시점 질의가 가능 & Zep invalidation\\
같은 종류끼리 관계는 자기 화살표 & READBACK\_OF, NEXT는 같은 노드 종류의 \emph{서로 다른} 두 노드 사이 관계. 조건을 간선에 적어 오해를 막음 & UML·ER·온톨로지 스키마 관례\\
\bottomrule
\end{tabularx}
\end{center}
{\footnotesize Airline의 \texttt{designator}는 ICAO 항공사 코드 3글자(Doc 8585, 예: DAL = Delta), Airport의 \texttt{location\_indicator}는 ICAO 공항 코드 4글자(Doc 7910, 예: KATL = 애틀랜타)다. \texttt{country}는 각각 항공사의 국적과 공항이 있는 나라다(RQ4).}

{\footnotesize ``지시 대상''은 그림에서 노드 종류 네 개(Runway, Taxiway, Ramp, Facility)를 묶은 이름이다. 노드 ID는 M3 규칙 \texttt{개체종류:공항코드:개체이름}을 따른다. KATL 목록에 없는 이름도 노드로 만들되 \texttt{in\_kb=false}로 표시한다.}


\begin{center}
\resizebox{0.85\textwidth}{!}{%
\begin{tikzpicture}[font=\small, x=0.13cm, y=1cm]
\draw[-{Stealth}, thick] (18,0) -- (102,0) node[right]{$t$ [s]};
\foreach \x in {20,30,...,100} { \draw (\x,0.08) -- (\x,-0.08) node[below,font=\scriptsize]{\x}; }
% transmissions (short labels; contents listed below the figure)
\fill[blue!25] (24.1,0.85) rectangle (30.4,1.15); \node[anchor=south,font=\scriptsize] at (27.2,1.15) {\#0005 ATCO};
\fill[orange!40] (31.7,1.85) rectangle (37.2,2.15); \node[anchor=west,font=\scriptsize] at (37.6,2.0) {\#0006 PILOT (복창)};
\fill[blue!25] (88.5,0.85) rectangle (91.2,1.15); \node[anchor=south,font=\scriptsize] at (89.8,1.15) {\#0022 ATCO};
% readback
\draw[red!60!black, thick, -{Stealth}] (33,1.83) to[out=-110,in=30] (29.5,1.5);
\node[font=\scriptsize, red!60!black, anchor=west] at (38,1.55) {READBACK\_OF \{MISMATCH\}};
\draw[red!60!black, thin] (37.9,1.55) -- (33.2,1.6);
% next
\draw[teal!60!black, thick, dashed, -{Stealth}] (30.4,0.75) to[out=-10,in=190] (88.5,0.75);
\node[font=\scriptsize, teal!60!black, fill=white, inner sep=1pt] at (72,0.3) {NEXT \{$\Delta t$ = 64.4 s\}};
% validity bars
\fill[green!45!black!60] (24.1,-0.95) rectangle (88.5,-0.7); \node[anchor=east,font=\scriptsize] at (23.5,-0.82) {$\tau$(TAXI\_VIA)};
\node[font=\scriptsize, anchor=west] at (88.8,-0.82) {MONITOR가 끝냄};
\fill[green!25] (24.1,-1.45) rectangle (30.4,-1.2); \node[anchor=east,font=\scriptsize] at (23.5,-1.33) {$\tau$(FOLLOW)};
\fill[green!45!black!60] (88.5,-1.95) rectangle (101,-1.7); \node[anchor=east,font=\scriptsize] at (88.0,-1.83) {$\tau$(MONITOR): 다음 지시가 없어 최대 1800 s까지};
% snapshot
\draw[dashed] (50,-2.2) -- (50,2.5) node[above,font=\scriptsize]{$\G|_{50}$};
\node[font=\scriptsize, anchor=north west] at (50.6,-2.15) {$\G|_{50}$: DAL1489는 TAXI\_VIA(M 경유 M2) 유효};
\end{tikzpicture}}
\end{center}
\noindent{\small \textbf{시간축 예시} (KATL 지상관제, 2024-08-26 00:00:00Z부터의 초). 파란 상자 = 관제사 송신, 주황 상자 = 조종사 송신, 초록 막대 = 지시의 유효 구간 $\tau$. 조종사 복창 \#0006은 \#0005가 끝난 뒤 1.3\,s 만에 시작한다(KATL 전체에서 연결된 복창 404건 중 351건이 5\,s 안). 50초 시점의 스냅샷 $\G|_{50}$에서는 TAXI\_VIA가 유효하고, 88.5초에 MONITOR가 이를 끝낸다. NEXT는 실제로 \#0005의 마지막 지시(TRAFFIC\_INFO)에서 \#0022 MONITOR로 이어진다.

\begin{tabularx}{\textwidth}{@{}ll>{\raggedright\arraybackslash}X@{}}
\#0005 & 24.1 s, ATCO & ``Delta fourteen eighty nine pass behind \dots\ nine left to Mike two taxi via Mike'' $\to$ FOLLOW, TAXI\_VIA (M 경유 M2), TRAFFIC\_INFO\\
\#0006 & 31.7 s, PILOT & ``Nine left Mike two via Mike \dots\ Delta fourteen eighty nine'' $\to$ TAXI\_VIA, FOLLOW (복창)\\
\#0022 & 88.5 s, ATCO & ``Delta fourteen eighty nine monitor twenty three eighty five'' $\to$ MONITOR\\
\end{tabularx}}
\end{document}
